\section{Details on approximating GP sample paths}
\label{sec:samplingAppendix}

In this section we give further details about the approach used in
Section~\ref{sec:sampling} to approximate a GP using random features. 
These random features can be used to approximate sample paths from the GP
posterior. By optimizing these sample paths we obtain
posterior samples over the global maxima $\xopt$.  We derive in more
detail the kernel approximation from (\ref{eq:kernel_approx}). Formally, the
theorem of \citep{Bochner:1959} states
\begin{theorem}[Bochner's theorem]
    A continuous, shift-invariant kernel is positive definite if and only if it
    is the Fourier transform of a non-negative, finite measure.
\end{theorem}
\noindent As a result given some kernel $k(\x,\x')=k(\x-\x',\zero)$ there must
exist an associated density $s(\vw)$, known as its \emph{spectral density},
which is the Fourier dual of $k$. This can be written as
\begin{align*}
    k(\vx,\vx')
    &= \int e^{-i\vw\T(\vx-\vx')} s(\vw)\,d\vw,
    \\
    s(\vw)
    &= \frac1{(2\pi)^d} \int e^{i\vw\T\vtau} k(\vtau,\zero) \,d\vtau.
\end{align*}
Further, we can treat this measure as a probability density
$p(\vw)=s(\vw)/\alpha$ where $\alpha=\int s(\vw)\,d\vw$ is the normalizing
constant. Consequently, the kernel can be written as
\begin{align*}
    k(\vx,\vx')
    &= \alpha \,\E_{p(\vw)}[e^{-i\vw\T(\vx-\vx')}]
    \\
    \intertext{and due to the symmetry of $p(\vw)$ \citep[see][]{Rasmussen:2006}
    we can write the expectation as}
    &= \alpha \,\E_{p(\vw)}[\tfrac12 (e^{-i\vw\T(\vx-\vx')} + e^{i\vw\T(\vx-\vx')}) ]
    \\
    &= \alpha \,\E_{p(\vw)}[\cos(\vw\T\vx-\vw\T\vx')]
    \\
    \intertext{We can then note that $\int_0^{2\pi}\cos(a+2b)\,db=0$ for any
    constant offset $a\in\R$. As a result, for $b$ uniformly distributed between
    0 and $2\pi$ we can write}
    &= \alpha \,\E_{p(\vw)}[\cos(\vw\T\vx-\vw\T\vx') +
                \E_{p(  b)}[\cos(\vw\T\vx+\vw\T\vx'+2b)]]
    \\
    &= \alpha \,\E_{p(\vw,b)}[\cos(\vw\T\vx+b - \vw\T\vx' - b) +
                              \cos(\vw\T\vx+b + \vw\T\vx' + b)]
    \\
    &= 2\alpha \,\E_{p(\vw,b)}[\cos(\vw\T\vx+b)\cos(\vw\T\vx' + b)]
    \\
    \intertext{The last equality can be derived from the sum of angles formula,
    which leads to the identity: $2\cos(x)\cos(y)=\cos(x-y)+\cos(x+y)$. Finally,
    we can average over $m$ weights and phases}
    &=
    \frac{2\alpha}{m} 
    \,\E_{p(\vW,\vb)}[\cos(\vW\vx+\vb)\T\cos(\vW\vx' + \vb)]
\end{align*}
where $[\vW]_i\sim p(\vw)$ and $[\vb]_i\sim p(b)$ are stacked versions of the
original random variables. The resulting quantity has the same expectation but
results in a lower variance estimator. If we let
$\vphi(\vx)=\sqrt{2\alpha/m}\cos(\vW\vx+\vb)$ denote a random $m$-dimensional
feature generated by this model we can also write the kernel as
$k(\x,\x')=\E_{p(\vphi)}[\vphi(\vx)\T\vphi(\vx')]$.

We now briefly show the equivalence between a Bayesian linear model using random
features $\vphi$ and a GP with kernel $k$. Consider now a linear model
$f(\vx)=\phi(\vx)\T\vtheta$ where $\vtheta\sim\Normal(\zero,\vI)$ has a standard
Gaussian distribution and observations $\D_n=\{(\vx_i,y_i)\}_{i\leq n}$ of the
form $y_i\sim\Normal(f(\vx_i), \sigma^2)$. The posterior of $\vtheta$ given
$(\D_n,\vphi)$ is also be normal $\Normal(\vm,\vV)$ where
\begin{align*}
    \vm &= (\vPhi\T\vPhi + \sigma^2\vI)^{-1}\vPhi\T\vy, \\
    \vV &= (\vPhi\T\vPhi + \sigma^2\vI)^{-1}\sigma^2,
\end{align*}
and where $[\vPhi]_i=\vphi(\vx_i)$ and $[\vy]_i=y_i$ consist of the stacked
features and observations respectively. We can also easily write the predictive
distribution over $f$ evaluated at a test point $\vx$, which is Gaussian
distributed with mean and variance given by
\begin{align*}
    \mu_n(\vx)
    &= \vphi(\vx)\T(\vPhi\T\vPhi + \sigma^2\vI)^{-1} \vPhi\T \vy,
    \\
    v_n(\vx)
    &= \vphi(\vx)\T(\vPhi\T\vPhi + \sigma^2\vI)^{-1} \vphi(\vx) \sigma^2.
\end{align*}
By a simple application of the matrix-inversion lemma these quantities can be
rewritten in terms which only make use of the inner products between features,
\begin{align*}
    \mu_n(\vx)
    &= \vphi(\vx)\T \vPhi\T(\vPhi\vPhi\T + \sigma^2\vI)^{-1}\vy_{1:t}
    \\
    v_n(\vx)
    &= \vphi(\vx)\T\vphi(\vx)
     - \vphi(\vx)\T\vPhi\T(\vPhi\vPhi\T + \sigma^2\vI)^{-1}\vPhi\vphi(\vx).
\end{align*}
the expectations of which are equivalent to the kernel $k$ and we obtain the
same expressions as that in (\ref{eq:gpvar}). 

